\documentclass[11pt]{article}
\usepackage[T1]{fontenc}
\usepackage{lmodern}
\usepackage{amsmath,amssymb,amsthm}
\usepackage[margin=1.05in]{geometry}
\usepackage[colorlinks=true,linkcolor=blue,citecolor=blue,urlcolor=blue]{hyperref}
\usepackage{microtype}
\newtheorem{theorem}{Theorem}[section]
\newtheorem{proposition}[theorem]{Proposition}
\newtheorem{corollary}[theorem]{Corollary}
\newtheorem{lemma}[theorem]{Lemma}
\theoremstyle{remark}
\newtheorem{remark}[theorem]{Remark}
\newcommand{\Q}{\mathbb{Q}}
\newcommand{\Z}{\mathbb{Z}}
\newcommand{\N}{\mathrm{N}}
\newcommand{\Irr}{\operatorname{Irr}}
\newcommand{\rad}{\operatorname{rad}}
\title{Irreducible compositions in an exceptional unicritical family}
\author{}
\date{}
\begin{document}
\begin{center}
{\Large Irreducible compositions in an exceptional unicritical family}
\end{center}
\vspace{0.6em}
\begin{abstract}
We construct irreducible compositions in a positive-shift subfamily left exceptional by a recent theorem on unicritically generated semigroups. Let $d\geq4$ be even and let $a\geq2$ be an integer which is not a $p$th power for any odd prime $p\mid d$, nor a fourth power when $4\mid d$. Put $f(x)=x^d+a-a^d$ and $g(x)=x^d+a$. We prove that $f\circ g\circ g\circ h$ is irreducible over $\Q$ for every composition $h$ of integral unicritical polynomials whose degrees have prime divisors among those of $d$. For a finite equal-degree generating set containing $f$ and $g$, this gives an explicit positive lower density. Within the exceptional four-constant alphabet, the shorter prefix $f\circ g$ suffices. The result includes infinitely many pairs in every degree divisible by four for which $g$ is reducible and $f$ fixes a square. The proof combines one real embedding, a field norm, and gaps between consecutive prime powers.
\end{abstract}

\section{Introduction and statements}
Let $\mathcal A$ be a finite set of polynomials of degree at least two in $\Z[x]$, and let $G=\langle\mathcal A\rangle$ be the semigroup they generate under composition. A natural question is whether the presence of one irreducible polynomial forces a positive proportion of the elements of $G$ to be irreducible over $\Q$, when polynomials are counted by degree. For generators of common degree $d$, write
\[
 \underline\delta(G)=\liminf_{B\to\infty}
 \frac{\#\{F\in G:\deg F\leq B,\ F\text{ irreducible over }\Q\}}
 {\#\{F\in G:\deg F\leq B\}}.
\]
The prime-degree case was treated in \cite{HJK}, following the quadratic work \cite{HJY}. For general common degree, Bhardwaj, Boyer-Paulet, Hindes, Qiu, and Sun \cite[Theorem~1.1]{BBHQS} establish the desired positive-proportion conclusion outside explicit exceptional families. In even degree, the exceptional constants are contained in
\begin{equation}\label{eq:exception}
 \{y^p-y^{pd},\ y^p,\ -y^p,\ -y^p-y^{pd}\},
 \qquad p\mid d\text{ prime},\quad y\in\Z.
\end{equation}
They expect the positive-proportion conclusion to extend to the exceptional cases. We treat a positive-shift subfamily of \eqref{eq:exception}.

Throughout, superscripts on integers denote ordinary powers. We write compositions explicitly, including $g\circ g$ rather than $g^2$, to distinguish iteration from multiplication. The empty composition is the identity polynomial.

\begin{theorem}\label{thm:main}
Let $d\geq4$ be even, and let $a\geq2$ be an integer satisfying
\begin{enumerate}
\item $a\notin\Z^p$ for every odd prime $p\mid d$;
\item if $4\mid d$, then $a\notin\Z^4$.
\end{enumerate}
Define
\[
 D=a^d-a,\qquad f(x)=x^d-D,\qquad g(x)=x^d+a,
 \qquad P=f\circ g\circ g.
\]
Then $P\circ h$ is irreducible over $\Q$ for every finite composition $h$ of polynomials $x^m+c$ with $c\in\Z$, $m\geq2$, and $\rad(m)\mid\rad(d)$. This includes the empty composition.
\end{theorem}

\begin{corollary}\label{cor:density}
Under the hypotheses of Theorem~\ref{thm:main}, let $G$ be generated by $r$ distinct polynomials $x^d+c_i$, with $c_i\in\Z$, including $f$ and $g$. Then
\[
 \underline\delta(G)\geq r^{-3}.
\]
If instead the generating set is a subset of
\begin{equation}\label{eq:alphabet}
 \{x^d+a-a^d,\ x^d+a,\ x^d-a,\ x^d-a-a^d\}
\end{equation}
containing $f$ and $g$, then $f\circ g\circ h$ is irreducible for every $h\in G$ and for the empty composition, and
\[
 \underline\delta(G)\geq r^{-2}.
\]
In particular, $\underline\delta(\langle f,g\rangle)\geq1/4$.
\end{corollary}

The main new instances occur when $4\mid d$ and $a=4z^4$. If $z$ is odd and nonzero, the hypotheses hold for every such $d$: the valuation $v_2(a)=2$ excludes every odd perfect power and every fourth power. For $d=2^k$, $k\geq2$, every nonzero integer $z$ is allowed. In these instances $a=(2z^2)^2$ and $f(a)=a$, while
\begin{equation}\label{eq:sophie}
 g(x)=\bigl(x^{d/2}-2zx^{d/4}+2z^2\bigr)
       \bigl(x^{d/2}+2zx^{d/4}+2z^2\bigr)
\end{equation}
is reducible. Thus an irreducible positive-shift generator is not available to obtain the conclusion by the existing norm argument. The first example is
\[
 f=x^4-252,\qquad g=x^4+4.
\]
The two-letter prefix $f\circ g$ gives the lower bound $1/4$ in the semigroup generated by this pair.

The argument in \cite[Proposition~3.8]{BBHQS} excludes the positive-shift partner in \eqref{eq:alphabet}. Its norm sufficient condition \cite[Proposition~3.1]{BBHQS} gives no conclusion at the first transition here, since $f(g(0))=a$ is a square in the displayed subfamily. Our proof deals with that transition using a negative real root of $f$ and then separates the fourth-power obstruction by its norm. All later transitions follow from elementary gaps between prime powers. We do not resolve the full exceptional locus \eqref{eq:exception}.

\section{Algebraic irreducibility criteria}
We use the classical binomial criterion: for a field $K$ of characteristic zero and $n\geq2$, the polynomial $x^n-b$ is irreducible over $K$ if and only if
\begin{equation}\label{eq:binomial}
 b\notin K^p\quad(p\mid n\text{ prime}),\qquad
 b\notin-4K^4\quad\text{if }4\mid n.
\end{equation}
See \cite[Chapter~VI, Section~9, Theorem~9.1]{Lang}. The exceptional fourth-power clause in \eqref{eq:binomial} is essential in composite even degree.

\begin{lemma}[Capelli's composition criterion]\label{lem:capelli}
Let $A\in K[x]$ be irreducible, let $\alpha$ be a root of $A$, and let $B\in K[x]$ be nonconstant. Then $A\circ B$ is irreducible over $K$ if and only if $B(x)-\alpha$ is irreducible over $K(\alpha)$.
\end{lemma}
\begin{proof}
If $\beta$ is a root of $B(x)-\alpha$, then $\alpha=B(\beta)$ and $K(\alpha)\subseteq K(\beta)$. The tower law gives
\[
 [K(\beta):K]=[K(\beta):K(\alpha)]\deg A.
\]
The right side equals $(\deg B)(\deg A)$ precisely when $B(x)-\alpha$ is irreducible over $K(\alpha)$. Since $\beta$ is a root of $A\circ B$, the equality is equivalent to irreducibility of that polynomial.
\end{proof}

The following familiar norm consequence is included to specify the parity and fourth-power issue exactly; compare \cite[Proposition~3.1]{BBHQS}.
\begin{lemma}\label{lem:norm}
Let $A\in\Q[x]$ be monic and irreducible of even degree $N$, and let $B=x^m+c$ with $m\geq2$ and $c\in\Q$. If $A(c)\notin\Q^p$ for every prime $p\mid m$, then $A\circ B$ is irreducible over $\Q$.
\end{lemma}
\begin{proof}
Suppose otherwise, and put $K=\Q(\alpha)$ for a root $\alpha$ of $A$. Lemma~\ref{lem:capelli} and \eqref{eq:binomial} imply either $\alpha-c=u^p$ for a prime $p\mid m$, or $4\mid m$ and $\alpha-c=-4u^4$, with $u\in K$. In the first case its norm is a rational $p$th power. In the second case its norm is
\[
 (-4)^N\N_{K/\Q}(u)^4,
\]
which is a rational square because $N$ is even. In both cases this contradicts the assumptions, since
\[
 \N_{K/\Q}(\alpha-c)=(-1)^N A(c)=A(c).
\]
\end{proof}

\section{The first mixed composition}
\begin{lemma}\label{lem:gap}
Let $d\geq4$ be even and $p\mid d$ be prime. If $b\geq2$ is an integer, then
\[
 (b^{d/p})^p-(b^{d/p}-1)^p>b^{d-d/p}\geq b^2.
\]
\end{lemma}
\begin{proof}
For an integer $u\geq2$, factor the difference of powers to obtain
\[
 u^p-(u-1)^p=\sum_{j=0}^{p-1}u^{p-1-j}(u-1)^j>u^{p-1}.
\]
Apply this with $u=b^{d/p}$. Since $d\geq4$ is even, $d-d/p\geq2$.
\end{proof}

\begin{proposition}\label{prop:first}
Under the hypotheses of Theorem~\ref{thm:main}, both $f$ and $f\circ g$ are irreducible over $\Q$.
\end{proposition}
\begin{proof}
For every prime $p\mid d$, Lemma~\ref{lem:gap} gives
\[
 (a^{d/p}-1)^p<a^d-a=D<(a^{d/p})^p.
\]
Thus $D\notin\Z^p$, hence $D\notin\Q^p$. Here we use the elementary fact that an integer which is a rational $p$th power is an integer $p$th power. Since $D>0$, it also does not belong to $-4\Q^4$. Criterion~\eqref{eq:binomial} proves $f$ irreducible.

Choose its negative real root $\alpha=-D^{1/d}$ and regard $K=\Q(\alpha)$ as a subfield of $\mathbb R$. The element $\alpha-a$ is negative in this embedding, so it is not a square in $K$. Also
\begin{equation}\label{eq:norma}
 \N_{K/\Q}(\alpha-a)=f(a)=a.
\end{equation}
For an odd prime $p\mid d$, equation~\eqref{eq:norma} and the hypothesis $a\notin\Z^p$ exclude $\alpha-a\in K^p$.

It remains to address the exceptional clause when $4\mid d$. If $\alpha-a=-4u^4$ for some $u\in K$, then, since $[K:\Q]=d$,
\[
 a=(-4)^d\N_{K/\Q}(u)^4
   =\bigl(4^{d/4}\N_{K/\Q}(u)\bigr)^4.
\]
This would make $a$ a rational fourth power, and therefore an integer fourth power, contrary to hypothesis. Thus \eqref{eq:binomial} proves $g(x)-\alpha=x^d-(\alpha-a)$ irreducible over $K$. Lemma~\ref{lem:capelli} proves $f\circ g$ irreducible over $\Q$.
\end{proof}

\section{Uniform propagation}
\begin{proof}[Proof of Theorem~\ref{thm:main}]
If $y\geq a^d+a$ is an integer, then $D<y$. For every prime $p\mid d$, Lemma~\ref{lem:gap} therefore gives
\begin{equation}\label{eq:valuegap}
 (y^{d/p}-1)^p<f(y)=y^d-D<(y^{d/p})^p.
\end{equation}
In particular, $f(y)$ is not a rational $p$th power for any prime $p\mid d$.

Put $H=f\circ g$, irreducible by Proposition~\ref{prop:first}. Since $g(a)=a^d+a$, equation~\eqref{eq:valuegap} applies to $H(a)$. Lemma~\ref{lem:norm} now proves $P=H\circ g$ irreducible.

For every $t\in\Z$, evenness of $d$ implies $g(t)=t^d+a\geq a$, and hence
\[
 g(g(t))\geq a^d+a.
\]
Equation~\eqref{eq:valuegap} yields
\begin{equation}\label{eq:allvalues}
 P(t)\notin\Q^p\qquad(t\in\Z,\ p\mid d\text{ prime}).
\end{equation}
We can therefore append the factors of $h$ one at a time. Suppose that $Q=P\circ h_0$ is already irreducible, and append $B=x^m+c$ as in the theorem. Then $h_0(c)\in\Z$, so
\[
 Q(c)=P(h_0(c))\notin\Q^p\qquad(p\mid m\text{ prime})
\]
by \eqref{eq:allvalues}. The polynomial $Q$ remains monic and has even degree. Lemma~\ref{lem:norm} proves $Q\circ B$ irreducible and completes the induction.
\end{proof}

\begin{proposition}\label{prop:short}
Under the hypotheses of Corollary~\ref{cor:density} for the alphabet \eqref{eq:alphabet}, $f\circ g\circ h$ is irreducible for every finite word $h$ in that alphabet, including the empty word.
\end{proposition}
\begin{proof}
Put $T=\{t\in\Z:|t|\geq a\}$. All four generator constants belong to $T$. The first three maps in \eqref{eq:alphabet} take $T$ into integers at least $a$, because $a^d-a\geq a$. For the last map, an input of magnitude $a$ maps to $-a$. For all other inputs in $T$, integrality gives $|t|\geq a+1$, and
\[
 (a+1)^d-a^d\geq d a^{d-1}>2a.
\]
Thus its output is greater than $a$. Consequently every word in the alphabet preserves $T$.

For $t\in T$ we have $g(t)\geq a^d+a$, so \eqref{eq:valuegap} implies that $H(t)=(f\circ g)(t)$ is not a rational $p$th power for any prime $p\mid d$. Starting from $H$ irreducible by Proposition~\ref{prop:first}, append letters as before. Each required norm value has the form $H(h_0(c))$, where the generator constant $c$ belongs to $T$ and $h_0(c)\in T$. Lemma~\ref{lem:norm} proves irreducibility at every step.
\end{proof}

\section{Counting by degree}
We include the elementary same-degree freeness argument to pass from words to actual polynomials; compare \cite[Proposition~3.9]{BBHQS}.
\begin{lemma}\label{lem:free}
In characteristic zero, distinct polynomials $x^d+c_1,\ldots,x^d+c_r$, with common degree $d\geq2$, generate a free semigroup under composition.
\end{lemma}
\begin{proof}
Degree determines word length. Suppose two words of equal length greater than one define the same polynomial. Removing their outer letters leaves monic polynomials $U,V$ of equal positive degree $M$, with
\[
 U^d+b=V^d+c.
\]
If $U\ne V$, then
\[
 \deg(U^d-V^d)=(d-1)M+\deg(U-V)>0,
\]
since $(U^d-V^d)/(U-V)$ has leading coefficient $d$ and degree $(d-1)M$. This contradicts $U^d-V^d=c-b$. Hence $U=V$ and $b=c$. Induction recovers all letters.
\end{proof}

\begin{proof}[Proof of Corollary~\ref{cor:density}]
By Lemma~\ref{lem:free}, there are exactly $r^n$ distinct polynomials of word length $n$, each of degree $d^n$. Every word of length $n\geq3$ with prefix $fgg$ is irreducible by Theorem~\ref{thm:main}; there are $r^{n-3}$ such words. Therefore, for $N=\lfloor\log_d B\rfloor\geq3$,
\[
 \frac{\#\{F\in G:\deg F\leq B,\ F\text{ irreducible}\}}
 {\#\{F\in G:\deg F\leq B\}}
 \geq\frac{\sum_{n=3}^{N}r^{n-3}}{\sum_{n=1}^{N}r^n}.
\]
Since $f\ne g$, we have $r\geq2$, and the right side tends to $r^{-3}$. For the restricted alphabet, Proposition~\ref{prop:short} supplies the prefix $fg$ of length two, giving the analogous limit $r^{-2}$.
\end{proof}

\begin{remark}
The restrictions on $a$ in Theorem~\ref{thm:main} are sufficient hypotheses for the real-place and norm proof; no necessity assertion is intended. The result does not settle exceptional pairs with the negative shift $x^d-a$ in place of $g$, when the positive-shift generator is absent. The mixed-degree assertion concerns irreducibility of the indicated compositions, not a density theorem for a mixed-degree semigroup.
\end{remark}

\begin{thebibliography}{9}
\bibitem{BBHQS}
A. Bhardwaj, A. Boyer-Paulet, W. Hindes, E. Qiu, and A. Sun,
\emph{Prime-powered images and irreducible polynomials in dynamical semigroups},
arXiv:2510.10310v1 (2025), \url{https://arxiv.org/abs/2510.10310}.
\bibitem{HJK}
W. Hindes, R. Jacobs, B. Keller, A. Kim, P. Ye, and A. Zhou,
\emph{On the proportion of irreducible polynomials in unicritically generated semigroups},
J. Algebra \textbf{677} (2025), 394--429.
\bibitem{HJY}
W. Hindes, R. Jacobs, and P. Ye,
\emph{Irreducible polynomials in quadratic semigroups},
J. Number Theory \textbf{248} (2023), 208--241.
\bibitem{Lang}
S. Lang, \emph{Algebra}, third revised edition,
Graduate Texts in Mathematics, vol.~211, Springer, New York, 2002.
\end{thebibliography}
\end{document}
